\subsection{Completion and full AA for arbitrary product alphabets}
\label{sec:resolution-product-alphabets}

This section removes a source-encoding obstruction. A source coordinate
may take values in an arbitrary finite alphabet, and one source query
returns its entire symbol. Both the quantum and classical algorithms
use this same convention. Alphabet sizes, atom probabilities, and the
number of source coordinates do not enter the bounds.

Let \(\Omega=\prod_{i=1}^n\Omega_i\), let
\(\mu=\bigotimes_i\mu_i\), and delete zero-mass symbols if necessary.
For Boolean \(G\) define its resampling influences by
\[
 I_i^\mu(G)=\Pr[G(X)\ne G(X_{-i},X_i')],
 \qquad I^\mu(G)=\sum_i I_i^\mu(G),
\]
where \(X\sim\mu\) and \(X_i'\sim\mu_i\) independently.
For \(x,y\in\Omega\) write
\(\Delta(x,y)=\{i:x_i\ne y_i\}\).
A fractional certificate for \(G\) at \(x\) is a vector
\(c^x\ge0\) with
\(\sum_{i\in\Delta(x,y)}c_i^x\ge1\) whenever \(G(y)\ne G(x)\).

\begin{lemma}[Product-alphabet certificate inequality]
\label{lem:resolution-alphabet-certificates}
Let \(U_b\subseteq G^{-1}(b)\) have mass \(\delta_b>0\), for
\(b=0,1\). Suppose every \(x\in U_b\) has a fractional certificate
of cost at most \(\gamma_b\). Then
\[
 \max_i I_i^\mu(G)\ge\frac{\delta_0\delta_1}{2\gamma_0\gamma_1}.
 \tag{AC}
\]
The same assertion holds after fixing any set of coordinates.
\end{lemma}

\begin{proof}
The only imported combinatorial ingredient is the finite Boolean-box
counting inequality proved as equation (4) in the proof of Lemma 4 of
\href{https://arxiv.org/abs/2611.12795}
{Huffstutler--Miloschewsky--Podder}.
We state exactly the form used. On a two-point product box \(C\), let
\(J\) be its free coordinates and let \(b_i(C)\) count its
unordered bichromatic \(i\)-edges. Consider some ordered antipodal
pairs \((x,y)\), with \(G(x)=1,G(y)=0\), and sets \(H_y\).
Suppose that replacing any subset of coordinates in \(H_y\cap J\)
of the paired \(x\) by their other box values leaves \(G=1\).
Suppose also that every \(z\in C\) with \(G(z)=1\) agreeing with
\(y\) on \(H_y\) satisfies \(w(\Delta(y,z))\ge1\).
Then the number \(N_C\) of these pairs obeys
\[
 N_C\le\sum_{i\in J}w_i b_i(C).                 \tag{BC}
\]
This is a statement about a finite Boolean cube, with no measure in it.
Relabeling each two-symbol coordinate by a bit preserves every
hypothesis. We now supply the measure-dependent extension.

Fix the stated certificates and define
\[
 w_i=2\gamma_0\,\mathbb E[c_i^X\mid X\in U_1],
 \qquad H_y=\{i:c_i^y>w_i\}\quad(y\in U_0).
\]
Then \(\sum_iw_i\le2\gamma_0\gamma_1\) and
\(w(H_y)\le\gamma_0\). Consequently at least half the conditional
\(\mu\)-mass of \(x\in U_1\) satisfies \(c^x(H_y)<1\).
Mark precisely these pairs \((x,y)\). Their product mass is at least
\(\delta_0\delta_1/2\). A marked pair cannot change the value of
\(G(x)\) by changing coordinates in \(H_y\), by its certificate.
If \(G(z)=1\) and \(z\) agrees with \(y\) on \(H_y\), the
certificate of \(y\) gives
\[
 1\le c^y(\Delta(y,z))\le w(\Delta(y,z)).
\]
Thus (BC) applies to the marked antipodal pairs in every box.

Take all boxes \(C=\prod_i C_i\) with \(|C_i|\in\{1,2\}\),
and assign the weight
\[
 a(C)=\prod_i a_i(C_i),\qquad
 a_i(\{u\})=\mu_i(u)^2,
 \quad a_i(\{u,v\})=\mu_i(u)\mu_i(v)\quad(u\ne v).
\]
Every ordered pair \((x,y)\) is antipodal in exactly one such box:
use \(C_i=\{x_i\}\) if the two symbols agree, and
\(C_i=\{x_i,y_i\}\) otherwise. Its box weight equals
\(\mu(x)\mu(y)\). Hence
\[
 \sum_C a(C)N_C\ge\delta_0\delta_1/2.             \tag{PAIR}
\]
For a fixed unordered \(i\)-edge with outside coordinates \(z\)
and endpoint symbols \(u,v\), summing the weights of its containing
boxes gives
\[
 \mu_i(u)\mu_i(v)
 \prod_{j\ne i}\left(\mu_j(z_j)^2+
               \sum_{a\ne z_j}\mu_j(z_j)\mu_j(a)\right)
 =\mu_i(u)\mu_i(v)\mu_{-i}(z).
\]
It follows that
\[
 \sum_C a(C)b_i(C)=\tfrac12 I_i^\mu(G).            \tag{EDGE}
\]
Multiply (BC) by \(a(C)\), sum, and use (PAIR), (EDGE), and the
bound on \(\sum_iw_i\). The result is
\(\delta_0\delta_1/2\le\gamma_0\gamma_1\max_i I_i^\mu(G)\),
which is (AC). Coordinate conditioning leaves a product measure, so
the final assertion follows from the same proof.
\end{proof}

\begin{theorem}[Completion on arbitrary product alphabets]
\label{thm:resolution-alphabet-completion}
Let \(f:D\to\{0,1\}\), where \(D\subseteq\Omega\), and suppose
there are nonnegative weights \(\lambda^x\), for \(x\in D\), with
\[
 \sum_i\lambda_i^x\le A,
 \qquad
 \sum_{i\in\Delta(x,y)}\min(\lambda_i^x,\lambda_i^y)\ge1
 \quad\text{if }f(x)\ne f(y).
 \tag{ADV}
\]
For every \(\eta>0\) there is a deterministic symbol-query tree of
depth at most \(\lceil128A^3/\eta^2\rceil\) whose unconditional
\(\mu\)-mass of errors on \(D\) is at most \(\eta\).
Constant partial functions need no queries.
\end{theorem}

\begin{proof}
Here are the product-measure details of the completion construction.
For \(a\in f^{-1}(1)\), let
\(d_a(z)=\sum_{i:z_i\ne a_i}\lambda_i^a\), and set
\(g(z)=\min\{1,\min_{a\in f^{-1}(1)}d_a(z)\}\).
Then \(g=0\) on the one-inputs and \(g=1\) on the zero-inputs:
for opposite inputs \(a,b\), (ADV) implies
\(d_a(z)+d_b(z)\ge1\) for every \(z\).
At any \(z\) with \(g(z)<1\), choose a minimizing \(a\).
Changing coordinate \(i\) to any other symbol can increase \(g\)
by at most \(\lambda_i^a\). Resampling symmetry therefore gives
\[
 \sum_i\mathbb E|g(X)-g(X_{-i},X_i')|\le2A.
\]
Integrating thresholds over \([1/4,3/4]\) provides a completion
\(G(z)=\mathbf1[g(z)<t]\) with \(I^\mu(G)\le4A\).
For every original domain point \(x\), the vector \(4\lambda^x\)
is a certificate for \(G\): if \(f(x)=1\), an opposite completion
value forces \(d_x(z)\ge t\); if \(f(x)=0\), a minimizing
one-input at \(z\) and the preceding two-distance inequality force
\(d_x(z)>1-t\). Both lower bounds are at least \(1/4\).

At a decision-tree transcript \(C\), let
\(\delta_b(C)=\mu(f^{-1}(b)\mid C)\) and
\(m(C)=\min_b\delta_b(C)\). The remaining measure is product,
and the restricted certificates have cost at most \(4A\).
Lemma~\ref{lem:resolution-alphabet-certificates} gives a free
coordinate with
\[
 I_i^{\mu\mid C}(G)\ge m(C)^2/(32A^2)
\]
unless \(m(C)=0\). Query that coordinate. The exact potential
identity is
\[
 I^{\mu\mid C}(G)
 -\mathbb E_{X_i\mid C} I^{\mu\mid C,X_i}(G)
 = I_i^{\mu\mid C}(G).
\]
Indeed, the other-coordinate resampling influences average unchanged.
After \(s\) queries output the majority domain value in the current
transcript. Its unconditional domain error is
\(e_s=\mathbb E[m(C_s)]\), and \(e_s\) is nonincreasing.
Telescoping the potential identity and using Jensen's inequality yields
\[
 s e_s^2\le\sum_{t<s}\mathbb E[m(C_t)^2]
 \le32A^2 I^\mu(G)\le128A^3.
\]
The stated depth suffices.
\end{proof}

\begin{lemma}[Quantum adversary weights do not depend on alphabet size]
\label{lem:resolution-symbol-adversary}
A \(t\)-query quantum algorithm deciding a partial Boolean function
with error at most \(1/3\) supplies weights (ADV) with \(A=O(t^2)\),
for arbitrary finite symbol alphabets.
\end{lemma}

\begin{proof}
Purify the algorithm. Let \(p_{r,i}(x)\) be its query probability
at index \(i\) immediately before query \(r\), and set
\(a_i(x)=\sum_{r=1}^t p_{r,i}(x)\); thus \(\sum_i a_i(x)=t\).
For opposite outputs, final states have overlap at most
\(2\sqrt2/3\). A single oracle call changes their overlap in
absolute value by at most
\(2\sum_{i:x_i\ne y_i}\sqrt{p_{r,i}(x)p_{r,i}(y)}\): identical
symbols act identically, and the difference of two unitary symbol
oracles has norm at most two. Telescoping and applying Cauchy--Schwarz
over \(r\) therefore gives
\[
 \sum_{i:x_i\ne y_i}\sqrt{a_i(x)a_i(y)}\ge
 c:=\tfrac12(1-2\sqrt2/3)>0.
\]
A second Cauchy--Schwarz inequality gives
\[
 c^2\le
 \left(\sum_{i:x_i\ne y_i}\min(a_i(x),a_i(y))\right)
 \left(\sum_{i:x_i\ne y_i}\max(a_i(x),a_i(y))\right)
 \le2t\sum_{i:x_i\ne y_i}\min(a_i(x),a_i(y)).
\]
Take \(\lambda_i^x=(2t/c^2)a_i(x)\).
\end{proof}

\begin{corollary}[Product-alphabet full AA with symbol queries]
\label{cor:resolution-product-alphabet-aa}
For every finite product measure on arbitrary finite alphabets,
every \(T\)-query quantum acceptance function \(Q\) has a
deterministic classical symbol-query simulator with mean error at
most \(e\) and \(O(T^6e^{-15})\) queries, for
\(T\ge1\) and \(0<e<1\). The absolute implied constant has no
dependence on alphabet sizes, biases, or number of coordinates.
\end{corollary}

\begin{proof}
Put \(m=\lceil16/e\rceil\), \(a=e/16\), and \(\eta=e/4\).
For each \(j=1,\ldots,m\), define a partial function by thresholding
\(Q\) at \(j/m\), with promise \(|Q-j/m|\ge a\).
Repeating the original quantum algorithm \(O(a^{-2})\) times and
thresholding the empirical acceptance frequency decides this partial
function with error at most \(1/3\). The preceding lemma and theorem
give a deterministic classifier \(D_j\) making
\(O(T^6a^{-12}\eta^{-2})\) symbol queries, whose unconditional
error mass on its promise is at most \(\eta\).

Output \(m^{-1}\sum_jD_j\). For each input, at most \(2am+2\)
thresholds lie within distance \(a\) of its acceptance value.
Rounding the integral of threshold indicators to this grid introduces
error at most \(1/m\). Therefore the mean approximation error is
at most
\[
 2a+3/m+\eta\le9e/16<e.
\]
Running the \(m\) classifiers costs \(O(T^6e^{-15})\) queries.
\end{proof}

The new ingredient relative to the Boolean completion theorem is the
weighted two-point-box summation. The symbol-oracle convention is
essential: this corollary does not claim that a whole large symbol can
be recovered from its bit encoding with one Boolean query. Its role
is as an internal source theorem in a transfer argument where both
forward coherent access and conditional source access are supplied
explicitly.
